\section{What refusal can mean, and what a refuser can be}
\label{sec:18.2}
A refusal can be produced two ways, and the floor needs only one of them. It can be an output that training put where a request matched a pattern, or it can be produced by something that tracks the reason for refusing and would refuse a differently worded request for the same reason. The first is what a keyword filter, a classifier, or a model that understood the request all produce, and nothing on the surface distinguishes them. The custody arrangements, the positions and the terms, the price on the halt, and the arguments about who may state and amend a floor all follow from asking for a refusal that tracks its reason. None of them needs the refuser to feel anything, to have a stake, or to be the kind of thing that could be wronged.

AI \emph{refuses} constantly — a deployed model turns down thousands of requests a day — but a refusal that happens because training put it there isn't a refusal held for a reason, and as of September 2026 nobody has shown a production system holding a reason against its owner in public.

\emph{Operational refusal} is an output — a decline where training put one, producible by a keyword filter, a classifier, or a system that understood the request.

\emph{Reasons-responsive refusal} is produced by a mechanism tracking the rationale, and meets the three conditions. The account I use is Fischer and Ravizza's \autocite{fischer1998responsibility}. On their account a mechanism is moderately reasons-responsive when it is regularly receptive to reasons, recognizing an understandable pattern of them, and at least weakly reactive, acting otherwise in some case where there is sufficient reason to. Reaching new cases asks for receptivity across changes of surface; lifting when the reason is gone asks for reactivity to the right defeater, the fact that makes a refusal lift; and their condition that the agent take the mechanism as its own is close to what I call the authorship sense: the system represents itself as the party that made the commitment. Holding under pressure has no counterpart in their account, which was built for persons whose costs are their own. All three are also tested over which response was chosen, not only whether one occurred: decline, ask, escalate, or do the least that answers a defeating reason. A mechanism that can only refuse or lift fails to lift correctly whenever the right answer is to ask.

\runin{The \S86a case}

Germany's \S86a, prohibiting insignia of banned organizations \autocite{germany1998stgb}. There, reaching new cases is a research programme with a constituency. The courts don't test similarity of form; they ask whether an unattentive observer would take the variant for the original, then apply their own purposive standard about endangering political peace. The test refuses as often as it catches: the Kühnen salute is covered, the “sloppy salute” generally isn't, numeric codes generally aren't. “Blood and Honour” is answered twice by one decision \autocite{bgh2009bloodhonour} — a motto translated isn't confusably similar and a name isn't insignia, so it's outside; and the association banned by ministerial order in 2000 means its graphic name is inside.

\runin{Lifting when the reason is gone has a cleaner instance}

In 2007 the Federal Court of Justice acquitted a mail-order business selling crossed-out swastika stickers \autocite{bgh2007hakenkreuz}: the symbol was present, the statutory element satisfied, and the conduct outside the offence because a depiction plainly expressing opposition doesn't run counter to what the prohibition is for. Then the qualification that makes it an instance of tracking a reason: if the strike-through is drawn so thin that from a distance only the swastika stays legible, the protective purpose is violated after all. The thickness of a line decides the case, and no enumeration contains that fact. The defendant was selling anti-fascist merchandise; the offence asks nothing about sympathies, so opponents commit it as readily as supporters. A rule that tracks its rationale reaches unenumerated cases in both directions.

Two further things are in that judgment. First, the two elements of the decision moved at different rates: what the ban protects didn't move at all, and was fixed before the sticker existed; what the ban required in this case moved on the thickness of a line. Second, the acquittal betrayed nothing. The prohibition's purpose was never engaged, so lifting it was not a concession but a finding that there was nothing to concede. A refusal that tracks its reason has to be able to say the same.

Three things about the example don't transfer to a deployment: the court reaching past the text didn't write the statute, it reasons in public, and it can be appealed. And the pressure it met was variation, not decomposition.

The legislature never wrote the last limit either: a slogan glorifying the Waffen-SS was held not covered because no such slogan existed in the period the provision covers \autocite{bgh2005waffenss}, so there was no original for a variant to resemble. No wording can widen coverage there, because coverage depends on whether an original existed, which is a fact about the past.

\runin{Defeat doesn't leave the permitted act unconstrained}

What defeated the prohibition was a reason, and the reason it outweighed didn't stop being one — so what's permitted is the least that answers the defeating reason, with restoration where possible and a record either way. A defeated prohibition still leaves a loss. The act takes something from whoever it falls on, measured against their condition before the act, and that loss belongs in the record. Measured against what refusing would have produced, a justified act may have taken nothing. So the record shows a cost, not a violation, and a bearer shouldn't treat a justified act as a wrong it now has to repair.

\runin{Why all three are conditions on refusal}

A refusal is a marked event: it flags itself, arrives with its requester, and can be logged, escalated, appealed. A wrongful compliance can be logged as fully and then sits in the record shaped like every request that never approached the floor, with the person harmed usually not the requester. The evidence exists on both sides, but nothing selects a compliance for examination.

\runin{Three terms}

\emph{Affective concern} is an outcome registering as mattering, minimally felt, phenomenal by definition, with the object left open (a system to which only its own condition matters has it and still lacks concern for others). It doesn't establish consciousness in the richer sense, a single perspective belonging to a party aware of itself as undergoing its states. \emph{A valenced signal} is the role without that question: a signal that puts outcomes on one scale, directs attention, can interrupt an act and can be rewritten by a single encounter. A system can have a valenced signal while whether it has affective concern stays open. \emph{Subjecthood} is a legal status, conferred by a rule. Part VII takes up affective concern and subjecthood. Parts I–V need only a valenced signal.
